\section{Datasets and Splits}
\label{sec:datasets_and_splits}

\subsection{Datasets}
\label{sec:datasets}

The pipeline developed in this work is based on four dataset sources, each with a specific role in the experimental protocol: (i) LOL-v2~\cite{DBLP:journals/tip/YangWHWL21}, (ii) MILLs~\cite{DBLP:journals/corr/abs-2511-15496}, (iii) ExDark~\cite{DBLP:journals/cviu/LohC19}, and (iv) LOL-v1~\cite{DBLP:conf/bmvc/WeiWY018}. All models are trained on a single source domain and are then evaluated without fine-tuning on a held-out in-domain test set and on independent external datasets. The external part of the protocol is intended to measure cross-dataset generalization, i.e., the ability of each model to enhance low-light images whose distribution differs from the training data. LOL-v1 is retained separately as a non-independent overlap-control set and is not included in the generalization interpretation.

The main training dataset is LOL-v2. This dataset contains both a real-captured subset, hereafter referred to as LOL-v2 Real, and a synthetic subset, referred to as LOL-v2 Synthetic. In the main protocol, only LOL-v2 Real is used, meaning that the low-light/normal-light pairs are acquired from real scenes rather than obtained through synthetic low-light degradation. LOL-v2 is a paired dataset: each low-light image is associated with a corresponding normally exposed reference image, which can be used as the target for supervised training and evaluation.

MILLs is used as a second paired cross-dataset evaluation set. In particular, the low-resolution digital single-lens reflex (DSLR) split, MILLs-DSLR, is adopted to keep the evaluation computationally manageable. This dataset introduces a controlled illumination setting. Its input images are associated with different illumination levels, and each input has a corresponding well-illuminated reference image. This makes it possible to evaluate the models also by studying how performance changes as the input illumination becomes more severe.

ExDark is used as an unpaired stress-test dataset. Here, unpaired means that the dataset provides low-light images without a corresponding normally exposed reference image. Therefore, ExDark is not used for reference-based comparison, but for no-reference evaluation and qualitative failure analysis on real-world dark scenes. Its category-based organization is preserved, meaning that images remain grouped according to the object category provided by the dataset, such as cars, people, animals, or indoor objects. This organization may help inspect whether some failure cases are related to the type of scene or object.

Finally, LOL-v1 is retained as a paired overlap-control set. A final exact-content audit showed that all 500 LOL-v1 low-light inputs match images contained in the LOL-v2 Real pool across its training, validation, and test partitions. It was nevertheless evaluated deliberately as a descriptive control, in order to inspect how the selected models behave on the earlier benchmark organization of the same underlying image content. LOL-v1 is therefore not interpreted as an independent cross-dataset benchmark and is excluded from all claims about generalization.

\subsection{Dataset Licenses and Terms of Use}

Since this project relies on publicly available datasets, their licenses and terms of use were reviewed to ensure that each dataset is used consistently with the conditions specified by its original authors.

For LOL-v2\footnote{\href{https://github.com/flyywh/CVPR-2020-Semi-Low-Light}{Official Source for LOL-v2}}, no explicit dataset license or terms of use were found in the official source consulted and linked in this report.

MILLs\footnote{\href{https://huggingface.co/datasets/davidserra9/MILL}{Official Source for MILLs}} is distributed through its Hugging Face dataset page, where the dataset metadata reports the license as MIT. This permissive license allows use, modification, and redistribution, provided that the copyright notice and permission notice are included in copies or substantial portions of the material.

ExDark\footnote{\href{https://github.com/cs-chan/Exclusively-Dark-Image-Dataset}{Official Source for ExDark}} is subject to dataset-specific usage restrictions. The repository reports the project as open source under the BSD-3 license and states that commercial use requires contacting the authors. This project uses ExDark exclusively within the scope of non-commercial academic research.

For LOL-v1\footnote{\href{https://daooshee.github.io/BMVC2018website/}{Official Source for LOL-v1}}, no explicit dataset license or terms of use were found in the official source consulted and linked in this report.

All datasets are used exclusively within the scope of this non-commercial academic work. Any applicable license terms or usage restrictions are respected, and the corresponding dataset papers are cited in the bibliography.

To make these conditions explicit, a dedicated dataset notice is provided both in the project repository and in the accompanying release. Any dataset files used or redistributed with this project remain subject to the original authors’ licenses and terms of use, and are provided solely to support non-commercial academic reproducibility.

\subsection{Split Construction}

Split construction follows the role assigned to each dataset, as discussed in Section~\ref{sec:datasets}. For the main source domain, LOL-v2 Real, the official training split is used to derive the training and validation partitions. In particular, a fixed fraction of the official training samples is deterministically held out as validation data using a fixed random seed. For this project, this held-out fraction is set to 10\% of the official training samples by the configuration file of the pipeline stage described in Section~\ref{sec:manifest_generation}. The official test split is kept untouched for the final in-domain evaluation. In this context, ``in-domain'' means that training, validation, and testing images come from the same dataset and acquisition setting.

The synthetic subset of LOL-v2 is kept separate from the main training protocol and is used only as a paired cross-dataset evaluation set. In this way, it provides a controlled test case for studying how models trained on real low-light acquisitions behave when evaluated on synthetically degraded low-light images.

For the independent external datasets, no additional train or validation partitions are created. LOL-v2 Synthetic, MILLs, and ExDark are each converted into a single evaluation partition, while preserving the metadata needed for grouped analysis. LOL-v1 is likewise stored as a single evaluation manifest, but it is handled separately in the interpretation because it is a deliberately non-independent overlap-control set. The resulting split sizes and roles are summarized in Table~\ref{tab:dataset_splits}.

\begin{table}[t]
	\centering
	\caption{Dataset split sizes used in the pipeline. The first row reports the main dataset, while LOL-v2 Synthetic, MILLs, and ExDark are independent external evaluation datasets. LOL-v1 (overlap) is retained as a non-independent descriptive control and is excluded from independent cross-dataset conclusions. Counts refer to samples, i.e., low-light input images. For paired datasets, each input sample is associated with a corresponding target image.}
	\label{tab:dataset_splits}
	\begin{tabular}{lccc}
		\toprule
		Dataset & Train & Validation & Test \\
		\midrule
		LOL-v2 Real & 620 & 69 & 100 \\
		\midrule
		LOL-v2 Synthetic & 810 & 90 & 100 \\
		MILLs & -- & -- & 499 \\
		ExDark & -- & -- & 7363 \\
		\midrule
		LOL-v1 (overlap) & -- & -- & 500 \\
		\bottomrule
	\end{tabular}
\end{table}